


\exo
Dans une classe, on connaît le nombre de filles, le nombre de garçons et l'effectif total.
\begin{enumerate}
\item Exprimer la proportion de filles dans la classe.
\item Exprimer la proportion de garçons dans la classe.
\end{enumerate}

\exo
Le prix initial du gaz est de 300 euros par an.
\begin{enumerate}
\item Calculer le prix après une augmentation de $10~\%$ la première année.
\item La seconde année, le prix augmente encore de $10~\%$. Calculer le nouveau prix.
\item Le prix a-t-il augmenté de $20~\%$ au total ? Justifier.
\end{enumerate}


\exo
Déterminer le pourcentage d'augmentation permettant de passer de $2700$ à $3672$.
On donnera la réponse arrondie à 0,1 \%.

\exo
Déterminer le pourcentage de diminution permettant de passer de $2300$ à $1541$.
On donnera la réponse arrondie à 0,1 \%.

\exo
Un employé a été payé 1471 \euro{} par mois durant l'année 2019.
Il était payé 1116 \euro{} par mois l'année précédente.
\begin{enumerate}
\item Calculer le taux d'évolution de son salaire entre les deux années.
On arrondira le résultat à 0,1 \% près.
\item Calculer la variation absolue de son salaire.
\end{enumerate}


\exo
Pour passer d'une quantité $Q_0$ à une quantité $Q_1$, on multiplie $Q_0$ par $1,18$.
Déterminer le taux d'évolution correspondant.

\exo
Après une augmentation de $4~\%$, un téléviseur coûte $213,20$ \euro{}.
Calculer son prix avant l'augmentation.

On donnera la réponse sous la forme d'un entier positif ou d'un nombre décimal arrondi au centième, suivi de l'unité qui convient.

\exo
Une quantité est multipliée par $2,5$. Déterminer le pourcentage d'augmentation correspondant.
On donnera la réponse arrondie à 0,1 \%.

\exo
Un minerai fournit $30~\%$ de sa masse en fonte.
Quelle quantité de minerai faut-il pour obtenir 840 kg de fonte ?

On donnera la réponse sous la forme d'un entier positif ou d'un nombre décimal arrondi au centième, suivi de l'unité qui convient.

\exo
Pour passer d'une quantité $Q_0$ à une quantité $Q_1$, on multiplie $Q_0$ par $0,95$.
Déterminer le taux d'évolution correspondant.
On donnera la réponse arrondie à 0,1 \%.


\exo
Après une augmentation de $8~\%$, le prix d'un objet est égal à $p$.
Quel était le prix initial ?


\exo
Un pantalon coûtait initialement $30$ \euro{}. Son prix augmente de $13~\%$ la première semaine, puis de $10~\%$ la semaine suivante.

Calculer le prix du pantalon après les deux augmentations.

On donnera la réponse sous la forme d'un entier positif ou d'un nombre décimal arrondi au centime près.

\exo
Le nombre d'adhérents à une association a augmenté de $5,9~\%$ en 2017, puis diminué de $1,9~\%$ en 2018.

Déterminer le taux d'évolution global correspondant à ces deux évolutions.

On arrondira le résultat à 0,1 \% près.

\exo
Les prix du gaz ont augmenté de $5~\%$ par an pendant six années successives.

Calculer le pourcentage d'augmentation global.
On arrondira le résultat à 0,01 \% près.

\exo
Le prix moyen de l'avocat par kilogramme était de $2,90$ \euro{} en janvier 1998.
Ce prix a diminué de $7~\%$ entre janvier 1998 et mai 1998, puis il a augmenté de $137~\%$ entre mai 1998 et avril 2016.

Quel était le prix moyen de l'avocat par kilogramme en avril 2016 ?

On donnera la réponse sous la forme d'un entier positif ou d'un nombre décimal arrondi au centime près.


\exo
On place $350$ \euro{} à $18~\%$ et $200$ \euro{} à $5~\%$.
Cela revient à placer $550$ \euro{} à un taux moyen de $x~\%$.
Combien vaut $x$ ?

On donnera la réponse arrondie à 0,1 \%.

\exo
Le chiffre d'affaires d'une entreprise a augmenté de $87~\%$ en sept ans.

Calculer le pourcentage d'évolution moyen annuel.
On arrondira le résultat à 0,01 \% près.
